\subsection{Sprint 0}

Na Sprint 0 era previsto que o grupo realizasse os mockups do Figma, as
User Stories e preparar toda a ambientação, populando wiki e criando fontes de
comunicação, entregando uma prévia da primeira versão do trabalho, além de
preparar tudo para começar o desenvolvimento. Quanto a minha parte fiquei de
ajudar nos mockups e estudar, principalmente através da participação em
workshops. 

O grupo foi capaz de concluir as tarefas propostas e eu também
consegui concluir as minhas. Entrando em mais detalhes, estudei um pouco de
WebDevelopment(HTML e CSS básico), além de relembrar como funcionava o uso
básico do FIGMA, por fim em cooperação com um colega de equipe AGES I, projetei
8 telas(6 que foram para o MediumFi e outras 2 que eram apenas alternativas),
sendo 2 de criação de pergunta para RH(checkbox e múltipla escolha), 2 de
visualização de pergunta para o Usuário(as mesmas) e upload de conteúdo externo
para RH + visualização de conteúdo externo para o Usuário.

Durante a realização dessas tarefas os problemas foram os que mais eram
esperados, afinal padronizações e burocracias são problemas bem comuns para
quem está trabalhando pela primeira vez em um projeto grande e uma equipe
grande, e assim foi para mim, deixar tudo padronizado e até ter medo de fazer a
primeira ação por não ter certeza de como o time trabalharia foram os maiores
problemas.

Por outro lado, achei interessante aprender mais a trabalhar em time e como
a criação de componentes/fontes/cores padronizados são essenciais para manter o
projeto bem estruturado e que usá-las ao próprio favor consegue poupar muito
tempo e deixar muito mais interligado o projeto.

Seguindo os aprendizados pretendo manter esta lógica de utilização de
componentes, reaproveitando e padronizando o máximo de código de meus outros
companheiros no FrontEnd para poder manter o estilo de programação uniforme,
ainda mais se utilizando dos plugins Prettier e ESLint, adotados pelo time para manter a uniformidade do código.